\documentclass[10pt,a4paper]{amsart}
\usepackage[margin=19mm]{geometry}
\usepackage{amsmath,amssymb,mathtools}
\usepackage[scr=rsfs,cal=euler]{mathalfa}
\usepackage{xfrac}
\usepackage[unicode,hidelinks,bookmarks=false]{hyperref}
\newcommand{\ie}{i.e.\ }
\newcommand{\cf}{cf.\ }
\newcommand{\resp}{respectively\ }
\newcommand{\Subsection}{\subsection}
\providecommand{\nobreakdash}{\nobreak\hbox{-}}
\setlength{\emergencystretch}{2em}
\multlinegap0pt
\usepackage{amssymb}
\usepackage{bm}
\usepackage{esint}
\newtheorem{theorem}{Theorem}
\newtheorem{lemma}{Lemma}
\newcommand{\ennorm}[1]{\lvert\!\lvert\!\lvert #1 \rvert\!\rvert\!\rvert}
\newcommand{\nc}{\mathrm{NC}}
\newcommand{\pw}{\mathrm{pw}}
\newcommand{\tri}{\mathcal{T}}
\title{A posteriori error estimates for a modified Morley FEM}
\author{A.K. Dond, D. Gallistl, S. Nayak, M. Schedensack}
\date{Source: arXiv:2602.14912v1 (CC BY 4.0)}
\begin{document}
\maketitle

\noindent\emph{Source and licence.} A posteriori error estimates for a modified Morley FEM. Version arXiv:2602.14912v1 (CC BY 4.0), distributed by arXiv under the Creative Commons Attribution 4.0 International licence, http://creativecommons.org/licenses/by/4.0/, as recorded in the arXiv licence metadata for this exact version and shown on the abstract page https://arxiv.org/abs/2602.14912v1. Full licence notice: \texttt{licenses/arxiv-2602.14912-CC-BY-4.0.txt}.\par\smallskip

\noindent\textit{Editorial scope.} Counted unit: the article's theorem t:sp\_reliability (source0.tex lines 509--518). The article's complete original proof of it is delivered with it (source0.tex lines 810--993, one \texttt{proof} environment of 184 lines, which the article places away from the statement). No mathematics is written, repaired or re-derived: the statement and the proof are the article's own lines, in the article's own notation, and no retained line is substituted or re-flowed.  Retained uncounted as the source-local context the proof needs: lines 360--369; lines 377--380; lines 465--468; lines 471--474; lines 533--540.  Nothing was omitted from any retained span, so the package carries no omission record.  No retained line contains a citation, so the wrapper carries no bibliography and no bibliography entry.  No retained line contains a figure environment, an image-inclusion command or drawing code, so no image is delivered and the package declares no asset dependency.  Editorial formatting only, in addition: an AMS \texttt{amsart} preamble that restates the article's own declarations and macros used by the retained lines, namely the environments \texttt{theorem}, \texttt{lemma} on separate counters, together with the standard packages those lines need.  The retained environments print in this document as Theorem 1, Lemma 1 in order of appearance, whereas the article prints the same items with the numbering of its own full document.  The complete native source of the article as distributed in the arXiv e-print for this exact version is bundled unchanged as \texttt{sources/194757/source0.tex}.
\begin{equation}\label{e:quasi_interpolation}
\begin{aligned}
    \sum_{j=0}^2 
    \lVert h_\tri^{-j}D_\pw^{2-j} (v-I_\mathrm{qi} v)\rVert_{L^2(\Omega)}
    &\lesssim \lVert D^2 v\rVert_{L^2(\Omega)},\\
    \lVert h_\tri^{-1}(v-I_\mathrm{qi} v)\rVert_{L^2(\Omega)}
    +\lVert \nabla_\pw (v-I_\mathrm{qi} v)\rVert_{L^2(\Omega)}
    &\lesssim \lVert \nabla v\rVert_{L^2(\Omega)}.
\end{aligned}
\end{equation}

\begin{align}\label{e:nodal_interpolation_approx}
    \sum_{j=0}^2 \lVert h_\tri^{-j}D^{2-j}(v-I_h v)\rVert_{L^p(T)} 
    \leq C(p) \lVert D^2 v\rVert_{L^p(T)}.
\end{align}

\begin{align}\label{e:sing_pert_weak_form}
    a_\varepsilon(u,v)=(f,v)_{L^2(\Omega)} 
    \qquad \text{for all }v\in H^2_0(\Omega),
\end{align}

\begin{align}\label{e:sing_pert_discrete_problem}
    a_{\varepsilon,h}(u_h,v_h) = (f,I_h v_h)_{L^2(\Omega)}
    \qquad \text{for all }v_h\in \mathcal{M}_0(\tri).
\end{align}

\begin{lemma}[a~posteriori estimate for nonconformity error]
    \label{l:sp_nc_estimator}
    Any $u_h\in\mathcal{M}_0(\tri)$ satisfies 
    \begin{align*}
        \inf_{v\in H^2_0(\Omega)} \ennorm{u_h-v}_{\varepsilon,\pw}
        \lesssim \mu_\nc + \mu_{I_h}.
    \end{align*}
\end{lemma}

\begin{theorem}[reliability]\label{t:sp_reliability}
    The exact solution $u\in H^2_0(\Omega)$ to \eqref{e:sing_pert_weak_form} 
    and the discrete solution $u_h\in \mathcal{M}_0(\tri)$ 
    to \eqref{e:sing_pert_discrete_problem} satisfy
    \begin{align*}
        \ennorm{u-u_h}_{\varepsilon,\pw}^2 
        + \lVert\nabla(u-I_h u_h)\rVert_{L^2(\Omega)}^2
        \lesssim \eta^2.
    \end{align*}
\end{theorem}

\begin{proof}[Proof of Theorem~\ref{t:sp_reliability}]
    \textbf{Step 1 (Error split).}
    The triangle inequality leads to
    \begin{align*}
        \lVert \nabla(u-I_h u_h)\rVert_{L^2(\Omega)}
        &\leq \lVert \nabla_\pw(u- u_h)\rVert_{L^2(\Omega)}
           + \lVert \nabla_\pw(u_h-I_h u_h)\rVert_{L^2(\Omega)}\\
        &\leq \ennorm{u-u_h}_{\varepsilon,\pw} + \mu_{I_h}.
    \end{align*}
    The error in the energy norm can be split as 
    \begin{align*}
        &\ennorm{u-u_h}_{\varepsilon,\pw}^2\\
        &\qquad
        =\inf_{v\in H^2_0(\Omega)} \ennorm{u_h-v}_{\varepsilon,\pw}^2
           + \sup_{v\in H^2_0(\Omega),\ennorm{v}_\varepsilon=1}
            \left( 
             (f,v)_{L^2(\Omega)} - a_{\varepsilon,\pw}(u_h,v) \right)^2.
    \end{align*}
    The first term is bounded with the help of 
    Lemma~\ref{l:sp_nc_estimator} by $\mu_\nc^2+\mu_{I_h}^2$.
    For the analysis of the second term,
    let $v\in H^2_0(\Omega)$ with $\ennorm{v}_\varepsilon=1$ and 
    set $e_v:=v-I_\mathrm{qi}v$ for the Morley quasi interpolation operator
    $I_\mathrm{qi}:H^2_0(\Omega)\to\mathcal{M}_0(\tri)$
    from \eqref{e:quasi_interpolation}.
    The expression inside the squared supremum then reads
    \begin{equation*}
    \begin{aligned}
        &(f,v)_{L^2(\Omega)} - a_{\varepsilon,\pw}(u_h,v)
        = \sum_{j=1}^4 T_j
    \end{aligned}
    \end{equation*}
    with 
    \begin{align*}
        T_1&:=(f,e_v)_{L^2(\Omega)} + (f, I_\mathrm{qi}v-I_hI_\mathrm{qi}v)_{L^2(\Omega)},
        &
        T_2&:=-\varepsilon^2 (D^2_\pw u_h,D^2_\pw e_v)_{L^2(\Omega)},\\
        T_3&:=- (\nabla I_h u_h,\nabla (v-I_h I_\mathrm{qi}v))_{L^2(\Omega)},
        &
        T_4&:= (\nabla (I_h u_h-u_h),\nabla v)_{L^2(\Omega)}.
    \end{align*}
    It remains to bound these terms by the error estimator.

    \textbf{Step~2 (Bound of $T_1$).}
    The estimates 
    \eqref{e:quasi_interpolation}
    of the nonconforming Morley quasi interpolation operator,
    the Cauchy inequality, and $\ennorm{v}_{\varepsilon}=1$ 
    prove
    \begin{align*}
        (f,v-I_\mathrm{qi}v)_{L^2(\Omega)} 
        &\lesssim 
          \min\{\lVert h_\tri f \rVert_{L^2(\Omega)} \,
            \lVert \nabla v \rVert_{L^2(\Omega)},
             \lVert h_\tri^2 f \rVert_{L^2(\Omega)} \,
            \lVert D^2 v \rVert_{L^2(\Omega)}\}\\
        &\leq \lVert\min\{h_\tri,h_\tri^2/\varepsilon\}
           f \rVert_{L^2(\Omega)}
         \leq \eta.
    \end{align*}
    Similarly, the $H^2$ approximation property of $I_h$ and
    the inverse estimate combined with the $H^1$ and $H^2$ stability of 
    $I_{qi}$ lead to
    $
        (f, I_\mathrm{qi}v-I_hI_\mathrm{qi}v)_{L^2(\Omega)}
        \lesssim \eta.
    $
    
    \textbf{Step~3 (Bound of $T_2$).}
    Since the jump and the average of 
    $D^2_\pw u_h\,\nu_F$ over $F$  is constant for any $F\in\mathcal{F}$
    and $[\nabla_\pw e_v]_F$ is affine on $F$ and vanishes in the midpoint of $F$,
    a piecewise integration by parts for $T_2$ leads to
    \begin{align}\label{e:proof_reliability_2}
        |T_2|
        =| \varepsilon^2 (D^2_\pw u_h,D^2_\pw e_v)_{L^2(\Omega)}|
        = \left |\varepsilon^2\sum_{F\in\mathcal{F}(\Omega)}
            ([D^2_\pw u_h\,\nu]_F,\nabla_\pw e_v)_{L^2(F)}
          \right|
          .
    \end{align}
    The multiplicative trace inequality proves for some
    $T_F\in\tri$ with $F\subseteq T_F$ that
    \begin{align}\label{e:multitrace}
        \|\nabla_\pw e_v\|_{L^2(F)}
        &\lesssim h_F^{-1/2} \|\nabla_\pw e_v\|_{L^2(T_F)}
            + \|\nabla_\pw e_v\|_{L^2(T_F)}^{1/2}
                \|D^2_\pw e_v\|_{L^2(T_F)}^{1/2}.
    \end{align}
    The approximation and stability properties from
    \eqref{e:quasi_interpolation} then prove
    \begin{align*}
    \|\nabla_\pw e_v\|_{L^2(F)}
        \lesssim \sqrt{\|\nabla  v\|_{L^2(\Omega_{T_F})} 
                    \,\|D^2  v\|_{L^2(\Omega_{T_F})}}
        \lesssim \varepsilon^{-1/2} \ennorm{v}_{\varepsilon,\Omega_{T_F}}.
    \end{align*}
    On the other hand, \eqref{e:multitrace}
    together with the approximation and stability properties
    from \eqref{e:quasi_interpolation} and Young's inequality 
    leads to
    \begin{align*}
        \|\nabla_\pw e_v\|_{L^2(F)}
        &\lesssim h_F^{-1/2} \lVert \nabla_\pw e_v\rVert_{L^2(T_F)}
           + h_F^{1/2} \lVert D^2_\pw e_v\rVert_{L^2(T_F)}\\
        &\lesssim h_F^{1/2} \lVert D^2 v\rVert_{L^2(\Omega_{T_F})}
        \lesssim (h_F^{1/2}/\varepsilon)\, \ennorm{v}_{\varepsilon,\Omega_{T_F}}.
    \end{align*}
    The combination of the foregoing displayed formulae leads to
    \begin{align}\label{e:proof_reliability_3}
        \|\nabla_\pw e_v\|_{L^2(F)}
        \lesssim \min\{\varepsilon^{-1/2},h_F^{1/2}/\varepsilon\}
        \ennorm{v}_{\varepsilon,\Omega_{T_F}}.
    \end{align}
    This and \eqref{e:proof_reliability_2} bound $T_2$ as
    \begin{align*}
        |T_2|
        \lesssim \sqrt{\sum_{F\in\mathcal{F}(\Omega)}
            \min\{\varepsilon^{3/2},h_F^{1/2}\varepsilon\}^2
               \lVert [D^2_\pw u_h\,\nu]_F\rVert_{L^2(F)}^2}.
    \end{align*}
    
    \textbf{Step~4 (Bound of $T_3$).}
    Since $\nabla I_h u_h$ is piecewise constant, a piecewise integration
    by parts and $v-I_hI_\mathrm{qi}v\in H^1_0(\Omega)$ 
    imply for $T_3$ that
    \begin{align*}
        -T_3=
        (\nabla I_h u_h,\nabla (v-I_h I_\mathrm{qi} v))_{L^2(\Omega)}
        = \sum_{F\in\mathcal{F}(\Omega)}
          ([\nabla I_h u_h\cdot\nu_F]_F, v-I_h I_\mathrm{qi} v)_{L^2(F)}.
    \end{align*}
    Note that the $H^2$ stability of the nodal interpolation 
    operator $I_h$ from \eqref{e:nodal_interpolation_approx},
    an inverse inequality, and the $H^1$ stability of $I_\mathrm{qi}$ 
    imply for any $T\in\tri$ that 
    \begin{align*}
        h_T^{-1}\lVert I_\mathrm{qi} v-I_h I_\mathrm{qi}v\rVert_{L^2(T)} 
        &+ \lVert \nabla(I_\mathrm{qi} v-I_h I_\mathrm{qi}v)\rVert_{L^2(T)}
        \lesssim h_T \lVert D^2 I_\mathrm{qi} v\rVert_{L^2(T)} \\
        &\qquad\qquad\qquad \qquad\qquad
        \lesssim \lVert \nabla I_\mathrm{qi} v\rVert_{L^2(T)} 
        \lesssim \lVert \nabla  v\rVert_{L^2(\Omega_T)},
    \end{align*}
    and, therefore, the operator $I_h I_\mathrm{qi}$ enjoys the 
    same approximation and stability properties 
    \eqref{e:quasi_interpolation} as $I_\mathrm{qi}$.
    Hence, the trace inequality and the 
    $H^1$ approximation and stability properties prove 
    \begin{align*}
        \lVert v-I_h I_\mathrm{qi}v\rVert_{L^2(F)}
        \lesssim h_T^{1/2}\lVert \nabla v\rVert_{L^2(\Omega_F)}
        \leq h_T^{1/2} \ennorm{v}_{\varepsilon, \Omega_F},
    \end{align*}
   while the $H^2$ approximation and stability properties show 
    \begin{align*}
        \lVert v-I_h I_\mathrm{qi}v\rVert_{L^2(F)}
        \lesssim h_T^{3/2}\lVert D^2 v\rVert_{L^2(\Omega_F)}
        \leq h_T^{3/2}\varepsilon^{-1} \ennorm{v}_{\varepsilon, \Omega_F}.
    \end{align*}
    The combination of these two estimates leads to 
    \begin{align*}
        \lVert v-I_h I_\mathrm{qi}v\rVert_{L^2(F)}
        \lesssim h_T^{1/2}\kappa_T \ennorm{v}_{\varepsilon, \Omega_F}.
    \end{align*}
    This and the finite overlap of the patches
    eventually show
    \begin{align*}
        |T_3|
        \lesssim  \sqrt{\sum_{F\in\mathcal{F}(\Omega)}
          h_T\kappa_T^2
          \,\lVert [\nabla I_h u_h\cdot\nu_F]_F\rVert_{L^2(F)}^2}
          .
    \end{align*}
    
    \textbf{Step~5 (Bound of $T_4$).}
    Finally, the last term $T_4$
    is estimated with a Cauchy inequality by
    \begin{align*}
       T_4= (\nabla (I_h u_h-u_h),\nabla  v)_{L^2(\Omega)}
        \leq \lVert \nabla (u_h-I_h u_h)\rVert_{L^2(\Omega)} .
    \end{align*}
    This concludes the proof.
\end{proof}

\end{document}
